\documentclass[11pt]{article}
\usepackage{lmodern}
\usepackage[T1]{fontenc}
\usepackage{amsmath,amssymb,amsthm}
\usepackage{geometry}
\usepackage[hidelinks]{hyperref}
\geometry{margin=1in}
\renewcommand{\labelitemi}{\(\bullet\)}

\newtheorem{theorem}{Theorem}
\newtheorem{lemma}[theorem]{Lemma}
\newtheorem{remark}[theorem]{Remark}
\newcommand{\Q}{\mathbb Q}
\newcommand{\F}{\mathbb F}
\newcommand{\Z}{\mathbb Z}
\newcommand{\Sha}{\mathop{\mathrm{Sha}}}
\newcommand{\Sel}{\operatorname{Sel}}
\newcommand{\Cl}{\operatorname{Cl}}
\newcommand{\Nm}{\operatorname{N}}
\newcommand{\CT}{\mathrm{CT}}

\title{An Explicit 2-Selmer and Cassels--Tate Criterion
for a Family of Quadratic Twists}
\author{\small\href{mailto:admin@cosmosgeek.org}{admin@cosmosgeek.org}}
\date{Working draft for specialist review}

\begin{document}
\maketitle

\begin{abstract}
We give an explicit cubic-field upper certificate and a
binary-quartic Cassels--Tate calculation for negative prime
twists of $E_0:y^2=x^3-28x+1$.
For the specified splitting conditions, the $2$-Selmer group
has dimension two and is generated by two explicit classes.
On the subfamily selected by a quadratic Frobenius sign, its
Cassels--Tate pairing is nondegenerate. This proves algebraic
rank zero and determines the full $2$-primary Tate--Shafarevich
group, without assuming its finiteness or any instance of BSD.
The Selmer-stability framework is due to Barrera Salazar,
Pacetti and Tornar\'ia; the binary-quartic pairing formula
is due to Fisher.
\end{abstract}

\section{Statement}

Put
\[
E_0:\ y^2=x^3-28x+1,\qquad
N=351124=4\cdot41\cdot2141.
\]
For a positive prime $q$, let
\[
E_q=E_0^{(-q)}:
y^2=x^3-28q^2x-q^3.
\]
Consider the integral binary quartics
\begin{align*}
G_1(s,t)&=-16s^4-24s^3t+24s^2t^2+8st^3-t^4,\\
G_2(s,t)&=-4s^4+16s^3t+12s^2t^2-28st^3+3t^4,\\
G_3(s,t)&=-4s^4+8s^3t+60s^2t^2+64st^3+15t^4.
\end{align*}

A prime $q$ is called \emph{eligible} if
\begin{equation}\label{eligible}
\begin{gathered}
q\nmid 2\cdot41\cdot2141,\qquad q\equiv7\pmod8,\\
\left(\frac{-q}{41}\right)
=\left(\frac{-q}{2141}\right)=1,\\
G_1(s,1)\bmod q
\text{ has exactly one root }r_q\in\F_q,
\text{ and that root is simple}.
\end{gathered}
\end{equation}
Equivalently, the last condition says that Frobenius has
cycle type $(1)(3)$ on the four roots of $G_1$.
Set
\[
\tau_q=\left(\frac{r_q(3r_q+1)}q\right).
\]
The argument below shows that the Legendre-symbol argument
is nonzero under \eqref{eligible}.

\begin{theorem}\label{main}
For every eligible prime $q$,
\[
\Sel_2(E_q/\Q)\simeq(\Z/2\Z)^2.
\]
Under the canonical identification of the $2$-torsion modules,
this group is the same subgroup of $H^1(\Q,E_0[2])$
for all eligible $q$.

The classes represented by $qz^2=G_2(s,t)$ and
$qz^2=G_3(s,t)$ form a basis, and their Cassels--Tate
pairing satisfies
\begin{equation}\label{pairing}
(-1)^{2\langle[G_2],[G_3]\rangle_{\CT}}
=\left(\frac{r_q(3r_q+1)}q\right)=\tau_q.
\end{equation}
Consequently, if $\tau_q=-1$, then
\[
\operatorname{rank}E_q(\Q)=0,\qquad
\Sha(E_q/\Q)[2^\infty]\simeq(\Z/2\Z)^2.
\]
\end{theorem}

\section{The cubic descent bound}

Let
\[
K=\Q(\theta),\qquad \theta^3-28\theta+1=0.
\]
Its defining polynomial has squarefree discriminant
\[
\operatorname{disc}(K)=87781=41\cdot2141,
\]
and therefore
\[
\mathcal O_K=\Z[\theta].
\]
The cubic is irreducible over $\Q$; in particular,
$E_0(\Q)[2]=E_q(\Q)[2]=0$.

\subsection{A finite class-group certificate}

We first establish
\begin{equation}\label{clbound}
\dim_{\F_2}\Cl(K)[2]\le1.
\end{equation}

The Minkowski bound is
\[
M_K=\frac{2}{9}\sqrt{87781}<65.84.
\]
Let $\mathfrak p_2=(2,\theta-1)$.
Exact norm calculation gives
\[
\Nm(12+3\theta-\theta^2)=-4,\qquad
\mathfrak p_2^2=(12+3\theta-\theta^2).
\]
The degree-two prime above $2$ is in the inverse ideal
class. The prime $3$ is inert, so its degree-three ideal
is $(3)$ and has norm $27$.

Here is a certificate for every remaining degree-one
prime ideal of norm below the Minkowski bound. An entry
$(p,r)$ denotes $(p,\theta-r)$.
\[
\begin{array}{rr|l|r}
p&r&\alpha&\Nm(\alpha)\\ \hline
7&3&7-4\theta-\theta^2&-14\\
7&5&5\theta-\theta^2&14\\
7&6&-2-11\theta-2\theta^2&7\\
13&1&\theta-\theta^2&26\\
17&8&2-5\theta-\theta^2&-34\\
29&26&-3\theta-\theta^2&-58\\
37&7&-16\theta-3\theta^2&37\\
41&16&-4+6\theta-\theta^2&-41\\
41&33&3+10\theta-2\theta^2&-41\\
43&23&7+4\theta-\theta^2&-86\\
47&2&2\theta-\theta^2&47\\
47&4&4\theta-\theta^2&47\\
47&41&-6\theta-\theta^2&47\\
61&33&56-16\theta-5\theta^2&-61
\end{array}
\]

The data are reproducible by substitution modulo $p$
and the norm identity
\begin{align*}
\Nm(a+b\theta+c\theta^2)
={}&a^3+56a^2c-28ab^2+3abc\\
&+784ac^2-b^3+28bc^2+c^3.
\end{align*}
Each displayed element vanishes at the indicated root;
its norm has absolute value $p$ or $2p$.
Thus each corresponding ideal is principal or equivalent
to $\mathfrak p_2$.

Reduction of $x^3-28x+1$ modulo the primes below the
Minkowski bound shows that this list, together with the
primes above $2$ and the inert prime above $3$,
accounts for all prime ideals of norm at most $65$.
The degree-two prime at $2$ is the only prime of residue degree two within that bound.

Every ideal class is therefore generated by
$[\mathfrak p_2]$, whose square is trivial.
This proves \eqref{clbound}.

\subsection{Even valuations and the real restriction}

Identify
\[
H^1(\Q,E_0[2])
\simeq
\ker\left(
K^\times/K^{\times2}
\xrightarrow{\Nm}
\Q^\times/\Q^{\times2}
\right).
\]
For $d=-q$, the descent polynomial is
\[
F_d(x)=x^3-28d^2x+d^3.
\]

We check the local integral hypotheses of
Barrera Salazar--Pacetti--Tornar\'ia
\cite[Definition 1.6 and Theorem 1.8]{BPT}.
For every finite prime $p\ne q$, the element $d$
is a $p$-adic unit, and the local cubic order generated
by $d\theta$ is maximal, since
$\operatorname{disc}(K)=41\cdot2141$ is squarefree.
This includes $p=2$.
At $p=q$, the cubic is an unramified field extension
of degree three, as proved in the next section.
Thus the local field alternative of Definition 1.6
holds there.

It follows that every global $2$-Selmer class is
represented by an $\alpha\in K^\times$ satisfying
\[
v_{\mathfrak p}(\alpha)\equiv0\pmod2
\quad\text{for every finite prime }\mathfrak p\text{ of }K.
\]
Let $V_{\rm even}$ be the space of these squareclasses.
There is an exact sequence
\[
0\longrightarrow
\mathcal O_K^\times/\mathcal O_K^{\times2}
\longrightarrow V_{\rm even}
\longrightarrow\Cl(K)[2]\longrightarrow0.
\]
As $K$ is totally real cubic,
\[
\dim_{\F_2}
\mathcal O_K^\times/\mathcal O_K^{\times2}=3.
\]
The class-group certificate consequently gives
\[
\dim_{\F_2}V_{\rm even}\le4.
\]

There are two independent constraints.

First, the norm must be a rational square.
Because all finite valuations are even, this amounts
to requiring its real sign to be positive.

Second, the local real Kummer image must have positive
sign at its distinguished embedding.
Write the roots of the defining cubic as
$\theta_1<\theta_2<\theta_3$.
Since $d=-q<0$, their order is reversed on $E_q$,
so the distinguished embedding is $\theta\mapsto\theta_3$.

These two restrictions are independent on units:
\[
\begin{array}{c|cc}
&\operatorname{sgn}\Nm&\operatorname{sgn}_{\theta_3}\\ \hline
-1&-&-\\
\theta&-&+
\end{array}
\]
because $\Nm(\theta)=-1$ and $\theta_3>0$.
Hence
\begin{equation}\label{upper}
\boxed{\dim_{\F_2}\Sel_2(E_q/\Q)\le4-2=2.}
\end{equation}
No finiteness assumption on $\Sha$ is involved.

\section{The two-dimensional Selmer plane}

All three quartics have classical invariants
\[
I=1344,\qquad J=1728,
\]
and polynomial discriminant
\[
\operatorname{disc}(G_i)=2^{12}\cdot41\cdot2141.
\]
Put $g_i=G_i/q$. Their invariants are
$I/q^2,J/q^3$.
Fisher's associated Jacobian
\[
Y^2=X^3-\frac{27I}{q^2}X-\frac{27J}{q^3}
\]
is identified with $E_q$ by
\[
x=\frac{q^2}{36}X,\qquad
y=\frac{q^3}{216}Y.
\]

Let $\varphi=12\theta/q$.
Fisher's cubic invariant
\[
z(g)=\frac{4a\varphi+3b^2-8ac}{3}
\]
for $g=as^4+bs^3t+cs^2t^2+dst^3+et^4$
gives
\begin{equation}\label{zinvariants}
\begin{aligned}
z(g_1)&=\frac{64}{q^2}(25-4\theta),\\
z(g_2)&=\frac{64}{q^2}(6-\theta),\\
z(g_3)&=\frac{64}{q^2}(11-\theta).
\end{aligned}
\end{equation}
Their norms are squares:
\[
\Nm(25-4\theta)=67^2,\quad
\Nm(6-\theta)=7^2,\quad
\Nm(11-\theta)=32^2.
\]
Moreover,
\begin{equation}\label{cubicproduct}
(25-4\theta)(6-\theta)(11-\theta)
=(40+9\theta-3\theta^2)^2.
\end{equation}
This gives the exact three-class relation
\[
[g_1]+[g_2]+[g_3]=0.
\]

\subsection{Local conditions are unchanged}

We verify the equality of local Kummer conditions
for two eligible twists $E_q,E_{q'}$.

At $2$, both $-q$ and $-q'$ are squares in $\Q_2$:
they are congruent to $1$ modulo $8$.

At $41$ and $2141$, both twist parameters are
local squares by \eqref{eligible}.
Thus the curves are locally isomorphic at these places.

At all other odd primes away from $q,q'$,
the twists have good reduction, and the local
$2$-Kummer image is the common unramified subgroup.
At infinity the twist parameters are both negative,
and the distinguished real embedding is the same.

It remains to consider $q$ and $q'$.

The polynomial
\[
P(u)=u^4+3u^3-6u^2-4u+1
\]
satisfies $P(2r)=-G_1(r,1)$.
Its cubic resolvent is
\[
R(Y)=Y^3+6Y^2-16Y-49,
\]
and
\[
R(X-2)=X^3-28X-1.
\]
Replacing $X$ by $-\theta$ identifies the resolvent
splitting field with that of $E_0[2]$.

A $(1)(3)$ Frobenius on the quartic induces a
$3$-cycle on the cubic roots.
Therefore $E_0[2]$ is unramified and irreducible
under Frobenius at an eligible prime $q$.

Writing $V=E_0[2]$, inflation--restriction gives
\[
H^1(\Q_q,V)=0:
\]
both $V/(\mathrm{Fr}_q-1)V$ and
the Frobenius invariants of
$\operatorname{Hom}(I_q,V)$ vanish.
Thus every local $2$-Kummer condition at $q$ is zero.

The same holds at $q'$.
We conclude that
\[
\boxed{\Sel_2(E_q/\Q)=\Sel_2(E_{q'}/\Q)
\subset H^1(\Q,V).}
\]
This is compatible with the published prime-twist
stability theorem \cite[Theorem 3.7]{BPT}.

\subsection{Local solubility and independence}

The preceding conditions also establish local solubility
of the three quartic coverings.
For completeness we give direct checks.

At $2,41,2141$, the twist parameter $-q$ is a
local square. Since
\[
G_1(1,0)=-16,\qquad
G_2(1,0)=G_3(1,0)=-4,
\]
all three equations $qz^2=G_i(s,t)$ have local points.

At the real place, use
\[
G_1(-1,1)=G_2(-1,1)=23>0,\qquad
G_3(0,1)=15>0.
\]

At $p=q$, the vanishing $H^1(\Q_q,V)=0$ gives
local solubility of every covering.

At every other odd prime
$p\nmid2q\cdot41\cdot2141$, each quartic has smooth
reduction and defines a genus-one curve over $\F_p$.
The Hasse bound gives a point over $\F_p$.
Smooth reduction and Hensel's lemma lift it to $\Q_p$.

Hence all three quartics are everywhere locally soluble
and represent elements of $\Sel_2(E_q/\Q)$.

To prove independence, use the splitting at $7$:
\[
x^3-28x+1
\equiv(x-3)(x-5)(x-6)\pmod7.
\]
The prime $7$ is not eligible, so $q\ne7$.
At the two factors $\theta=3,5$, the classes
$6-\theta$ and $11-\theta$ have quadratic signs
\[
\begin{array}{c|cc}
&\theta=3&\theta=5\\ \hline
6-\theta&-&+\\
11-\theta&+&-
\end{array}
\]
in $\F_7^\times/\F_7^{\times2}$.
They are independent.

Together with \eqref{upper}, this proves
\begin{equation}\label{selmerplane}
\boxed{
\Sel_2(E_q/\Q)
=\langle[g_2],[g_3]\rangle
\simeq\F_2^2.
}
\end{equation}

\section{The explicit Cassels--Tate pairing}

We apply Fisher's binary-quartic formula
\cite[Theorem 3.1 and Remark 3.2]{Fisher}.

The product identity \eqref{cubicproduct} gives the
required exact normalization. In Fisher's notation take
\[
m=\frac{512}{q^3}(40+9\theta-3\theta^2),
\]
so that
\[
z(g_1)z(g_2)z(g_3)=m^2.
\]
The quadratic covariant from his equation (5),
expanded in $1,\varphi,\varphi^2$, has
$\varphi^2$-coefficient
\begin{equation}\label{covariant}
\gamma_1(s,t)=-\frac{32}{9q}s(3s+t).
\end{equation}
For an explicit check, at $q=1$ the coefficients of $s^2,st,t^2$ in Fisher's $H_1$ in the $\theta$ basis are respectively
\[
64(25-4\theta),\quad 128-192\theta,\quad 448+64\theta-32\theta^2.
\]
Multiply these by $8(6-\theta)(11-\theta)/(40+9\theta-3\theta^2)$, reduce modulo $\theta^3-28\theta+1$ and convert $\theta^2$ to $\varphi^2/144$. For arbitrary $q$ the products scale by $q^{-3}$, while $\theta^2=q^2\varphi^2/144$, giving \eqref{covariant}.
Also
\[
g_2(1,0)=-\frac4q.
\]

The first Hilbert argument has squareclass $-q$,
and
\[
\gamma_1(s,t)=
(-q)\,2s(3s+t)
\left(\frac4{3q}\right)^2.
\]

Fisher's formula therefore initially gives
\[
\prod_v(-q,-q\cdot2s_v(3s_v+t_v))_v.
\]
The constant factor $(-q,-q)_v$ must not be
discarded place by place. It may be removed
from the \emph{global product} because Hilbert
reciprocity gives
\[
\prod_v(-q,-q)_v=1.
\]
Thus
\begin{equation}\label{hilbertproduct}
(-1)^{2\langle[g_1],[g_2]\rangle_{\CT}}
=\prod_v(-q,2s_v(3s_v+t_v))_v,
\end{equation}
where $qz_v^2=G_1(s_v,t_v)$ and the second
Hilbert argument is nonzero.

By \eqref{cubicproduct} and alternation,
\[
\langle[g_1],[g_2]\rangle_{\CT}
=\langle[g_2],[g_3]\rangle_{\CT}.
\]

\subsection{Evaluation away from the twisting prime}

At $2,41,2141$, the first Hilbert argument
$-q$ is a local square. Every factor is $+1$.
The choice $(s,t)=(1,0)$ has
$2s(3s+t)=6\ne0$.

At infinity, choose $(-1,1)$:
\[
G_1(-1,1)=23>0,\qquad
2s(3s+t)=4>0.
\]
The real Hilbert symbol is $+1$.

At an odd prime $p\ne q$, if both Hilbert
arguments are units, the Hilbert symbol is $+1$.
We therefore need local points with the second
argument nonzero, preferably a unit.

For $p\ge17$ away from the discriminant,
the smooth genus-one reduction has at least
\[
p+1-2\sqrt p>8
\]
points.
At most eight lie over the four branch
directions or the two directions
$s(3s+t)=0$.
Hence an admissible point exists with
$G_1/q$ a nonzero square and
$2s(3s+t)$ a unit.
Hensel lifting gives the desired local point.

For the remaining small primes, the following
finite checks suffice.
Write $h=2s(3s+t)$.
\[
\begin{array}{c|cc}
p&
(s,t)\text{ if }(q/p)=1&
(s,t)\text{ if }(q/p)=-1\\ \hline
5&(-1,-1)&(-1,1)\\
7&(-1,1)&(-1,-1)\\
11&(-1,1)&(-1,-1)\\
13&(-1,-1)&(-2,1)
\end{array}
\]
The corresponding values of $G_1$ modulo $p$
have the indicated quadratic character, and
$h$ is nonzero modulo $p$.
Thus $G_1/q$ is a nonzero local square.

At $p=3$, use $(1,10)$ when $q\equiv1\pmod3$:
\[
G_1(1,10)=144=12^2,\qquad h=26.
\]
When $q\equiv2\pmod3$, use $(-1,1)$:
\[
G_1(-1,1)=23,\qquad h=4.
\]
These choices also give local squares and unit
second arguments.

Thus every place $v\ne q$ contributes $+1$
to \eqref{hilbertproduct}.

\subsection{The twisting-prime contribution}

At $q$, let $r_q$ be the unique simple root
of $G_1(s,1)$ modulo $q$.

The derivative is nonzero, so an integral lift
$s_q\equiv r_q\pmod q$ can be chosen such that
\[
G_1(s_q,1)/q
\]
is a nonzero square in $\Q_q$.
Indeed, varying $s_q$ by a multiple of $q$
varies the quotient modulo $q$ through every
residue class.

The second argument is a unit.
We have $r_q\ne0$ because $G_1(0,1)=-1$.
Also $3r_q+1\ne0$, since
\[
G_1(-1/3,1)=-25/81
\]
and eligible $q$ is neither $3$ nor $5$.

Therefore
\[
(-q,2s_q(3s_q+1))_q
=\left(\frac{2r_q(3r_q+1)}q\right).
\]
Since $q\equiv7\pmod8$,
\[
\left(\frac2q\right)=1,
\]
and hence
\[
\prod_v(-q,2s_v(3s_v+t_v))_v
=\left(\frac{r_q(3r_q+1)}q\right)=\tau_q.
\]
This proves \eqref{pairing}, including its sign
orientation and all local factors.

\section{Rank and Tate--Shafarevich group}

Assume $\tau_q=-1$.
The pairing is then nonzero on the basis
$[g_2],[g_3]$ of the two-dimensional
$2$-Selmer group.

An alternating pairing on a two-dimensional
$\F_2$-space is nondegenerate if it is nonzero.
The rational Kummer image lies in the radical
of the Cassels--Tate pairing. Consequently
\[
E_q(\Q)/2E_q(\Q)=0.
\]
Since $E_q(\Q)[2]=0$, the Mordell--Weil theorem gives
\[
\operatorname{rank}E_q(\Q)=0.
\]

The exact Selmer sequence now gives
\[
\Sha(E_q/\Q)[2]\simeq\F_2^2.
\]

We can also determine the entire $2$-primary group,
without assuming finiteness.
Suppose $\xi\in\Sha(E_q/\Q)$ has order four.
Then $2\xi$ is a nonzero element of $\Sha[2]$.
For every $\eta\in\Sha[2]$,
bilinearity gives
\[
\langle2\xi,\eta\rangle_{\CT}
=\langle\xi,2\eta\rangle_{\CT}=0.
\]
This contradicts nondegeneracy on $\Sha[2]$.
Thus $\Sha[4]$ contains no element of order four.
Any element of higher $2$-power order would
produce one of order four, so none exists.
Therefore
\[
\boxed{\Sha(E_q/\Q)[2^\infty]\simeq(\Z/2\Z)^2.}
\]
This completes the proof of Theorem \ref{main}.

\section{Nonemptiness and Chebotarev}

The seed $q=31$ satisfies all eligibility conditions.
Modulo $31$, the quartic $G_1(s,1)$ has the
unique simple root $r_{31}=3$.
Moreover
\[
31\equiv7\pmod8,\qquad
\left(\frac{-31}{41}\right)
=\left(\frac{-31}{2141}\right)=1,
\]
and
\[
\tau_{31}
=\left(\frac{3(3\cdot3+1)}{31}\right)
=\left(\frac{30}{31}\right)=-1.
\]

For clarity, the local pairing sign at the seed
can also be checked without symbolic reduction:
\[
G_1(3,1)=-1705=31(-55).
\]
Since $-55$ is a square in $\Q_{31}$,
this is an admissible local point.
The second Hilbert argument is $60$, giving
\[
(-31,60)_{31}
=\left(\frac{29}{31}\right)=-1.
\]

To see that the negative-sign subfamily has
positive prime density, let $M$ be the compositum of:

\begin{itemize}
\item the splitting field of
$P(u)=u^4+3u^3-6u^2-4u+1$;
\item the extensions obtained by adjoining the
square roots of $r(3r+1)$ for its four corresponding
quartic roots;
\item the cyclotomic field
$\Q(\zeta_{8\cdot41\cdot2141})$.
\end{itemize}

The discriminant of the monic polynomial $P$ is $41\cdot2141$.
The resultant of $G_1(r,1)$ and $r(3r+1)$ is $25$.
Thus these defining data have no ramification
at $31$; in particular the Frobenius class of
$31$ in $\operatorname{Gal}(M/\Q)$ is defined.

Every prime whose Frobenius belongs to this
conjugacy class has the same quartic cycle type,
the required congruence and Legendre-symbol
conditions, and the same quadratic sign at its
unique fixed root. These conditions are
conjugacy-invariant.

Chebotarev therefore supplies a positive-density
set of eligible primes with $\tau_q=-1$.
The theorem applies to every prime in that set.

\section{Scope}

This is an arithmetic statement about a specified
quadratic-twist family. It proves algebraic rank
zero and determines the full $2$-primary
Tate--Shafarevich group on the negative-pairing
subfamily.

It does not assert analytic rank zero, the
leading-coefficient formula of BSD, or BSD for
arbitrary elliptic curves.

The general prime-twist Selmer stability is part
of \cite{BPT}, and the binary-quartic
Cassels--Tate formula is part of \cite{Fisher}.
The explicit cubic certificate and local Hilbert
product above are the calculations needed for
this particular family. No claim of priority or
literature-wide originality is made.

\begin{thebibliography}{9}

\bibitem{BPT}
D. Barrera Salazar, A. Pacetti and G. Tornar\'ia,
\emph{On $2$-Selmer groups and quadratic twists
of elliptic curves},
\href{https://arxiv.org/abs/2001.02263}{arXiv:2001.02263}, version 3 (2020).
In particular Definition 1.6, Theorem 1.8,
Theorem 2.16 and Theorem 3.7.

\bibitem{Fisher}
T. Fisher,
\emph{On binary quartics and the Cassels--Tate pairing},
Research in Number Theory \textbf{8} (2022),
article 74.
\href{https://doi.org/10.1007/s40993-022-00376-z}{doi:10.1007/s40993-022-00376-z}.
In particular Theorem 3.1 and Remark 3.2.

\end{thebibliography}

\end{document}
